\documentclass[10pt,a4paper]{amsart}
\usepackage[margin=19mm]{geometry}
\usepackage{amsmath,amssymb,mathtools}
\usepackage[scr=rsfs,cal=euler]{mathalfa}
\usepackage{xfrac}
\usepackage[unicode,hidelinks,bookmarks=false]{hyperref}
\newcommand{\ie}{i.e.\ }
\newcommand{\cf}{cf.\ }
\newcommand{\resp}{respectively\ }
\newcommand{\Subsection}{\subsection}
\providecommand{\nobreakdash}{\nobreak\hbox{-}}
\setlength{\emergencystretch}{2em}
\multlinegap0pt
\usepackage{amssymb}
\usepackage{enumerate}
\usepackage{xcolor}
\newtheorem{lemma}{Lemma}
\newcommand{\pc}{p_{\mathsf{c}}}
\title{The $q<1$ Random-Cluster Model on Wired Trees: Uniqueness and Negative Dependence}
\author{Heehyun Park}
\date{Source: arXiv:2608.08565v1 (CC BY 4.0)}
\begin{document}
\maketitle

\noindent\emph{Source and licence.} The $q<1$ Random-Cluster Model on Wired Trees: Uniqueness and Negative Dependence. Version arXiv:2608.08565v1 (CC BY 4.0), distributed by arXiv under the Creative Commons Attribution 4.0 International licence, http://creativecommons.org/licenses/by/4.0/, as recorded in the arXiv licence metadata for this exact version and shown on the abstract page https://arxiv.org/abs/2608.08565v1. Full licence notice: \texttt{licenses/arxiv-2608.08565-CC-BY-4.0.txt}.\par\smallskip

\noindent\textit{Editorial scope.} Counted unit: the article's lemma lem:uniform-branch-survival (source0.tex lines 1701--1704). The article's complete original proof of it is delivered with it (source0.tex lines 1706--1887, one \texttt{proof} environment of 182 lines). No mathematics is written, repaired or re-derived: the statement and the proof are the article's own lines, in the article's own notation, and no retained line is substituted or re-flowed.  Retained uncounted as the source-local context the proof needs: lines 1647--1668.  Nothing was omitted from any retained span, so the package carries no omission record.  No retained line contains a citation, so the wrapper carries no bibliography and no bibliography entry.  No retained line contains a figure environment, an image-inclusion command or drawing code, so no image is delivered and the package declares no asset dependency.  Editorial formatting only, in addition: an AMS \texttt{amsart} preamble that restates the article's own declarations and macros used by the retained lines, namely the environments \texttt{lemma} on separate counters, together with the standard packages those lines need.  The retained environments print in this document as Lemma 1 in order of appearance, whereas the article prints the same items with the numbering of its own full document.  The complete native source of the article as distributed in the arXiv e-print for this exact version is bundled unchanged as \texttt{sources/207173/source0.tex}.
\begin{lemma}
\label{lem:conditional-star-factorization}
Let $0<q<1$, and let $\nu$ be an infinite-volume measure satisfying the wired-tree DLR condition. Fix a vertex $u$ with children $u_1,\ldots,u_d$, and write $e_i:=\{u,u_i\}$. Let $\mathcal H_u := \sigma\left(\omega(e):e\notin\{e_1,\ldots,e_d\}\right)$. Thus $\mathcal H_u$ fixes every edge except the $d$ edges from $u$ to
its children. On $\mathcal R_u^c$, their conditional distribution given
$\mathcal H_u$ is the product Bernoulli distribution with parameter
$\widehat p$. More precisely, for every
$(a_1,\ldots,a_d)\in\{0,1\}^d$,
\[
    \mathbf 1_{\mathcal R_u^c}
    \nu\left(
        \omega(e_i)=a_i,\ i\in[d]
        \,\middle|\,
        \mathcal H_u
    \right)
    =
    \mathbf 1_{\mathcal R_u^c}
    \prod_{i=1}^d
    \widehat p^{a_i}
    (1-\widehat p)^{1-a_i}
\]
almost surely.
\end{lemma}

\begin{lemma}[Uniform branch survival]
\label{lem:uniform-branch-survival}
Let $0<q<1$, $\Delta\ge3$, and $p>\pc$. There exists $\theta\in(0,1]$, depending only on $p,q,$ and $d$, such that, for every infinite-volume measure $\nu$ satisfying the wired-tree DLR condition and every vertex $u$, $\nu\left(\mathcal B_u\,\middle|\,\mathcal F_u^{\mathrm{ext}}\right)\ge\theta $ almost surely.
\end{lemma}

\begin{proof}
Choose $\varepsilon\in(0,1)$ sufficiently small so that $d(1-\varepsilon)\widehat p>1$. For $b\in\{0,\ldots,d\}$, define
\[
    \Psi_b(z)
    :=
    1-(1-\varepsilon)^b
    \left(
        1-(1-\varepsilon)\widehat p z
    \right)^{d-b}.
\]
We have $\Psi_0(0)=0$ and $\Psi_0'(0)=d(1-\varepsilon)\widehat p>1$, whereas $\Psi_b(0)=1-(1-\varepsilon)^b>0$ for $b\ge1$. Hence there exists $\delta\in(0,1)$, depending only on $p,q,$ and $d$, such that $\Psi_b(\delta)\ge\delta$, for $b=0,\ldots,d$.

Independently of $\omega\sim\nu$, assign to every oriented child edge $e$ a Bernoulli mark $I_e$ with parameter $\varepsilon$, and denote their joint law by
\[
    \overline\nu
    :=
    \nu
    \otimes
    \bigotimes_e \operatorname{Ber}(\varepsilon),
\]
where the product runs over all oriented child edges. Edges with $I_e=1$ are called \emph{probes}. Starting from $u$, run the following breadth-first exploration. At each active vertex $v$, inspect the events $\mathcal B_w$ for its probe children $w$, without inspecting the corresponding probe edges. If a probe child is blue, select the first such probe edge and stop. For every nonprobe child edge, reveal its state and declare its child active when the edge is open.

We first show that, conditional on $\mathcal B_u^c$, with probability at least $\delta$ the exploration selects a probe edge $e=\{v,w\}$ whose parent $v$ is active and whose child $w$ is blue. This is the unrevealed contact edge whose conditional open and closed probabilities will be compared below.

The star calculation in Lemma~\ref{lem:conditional-star-factorization} gives the following slightly more general fact. Fix an active vertex $v$ with children $v_1,\ldots,v_d$, write $e_i:=\{v,v_i\}$, and condition on the configuration $\eta$ off these $d$ edges. Let $S:=\{i:\mathcal B_{v_i}\text{ occurs}\}$ and suppose $|S|=b$. On $\mathcal B_v^c$, every edge $e_i$ with $i\in S$ must be closed. For $A\subseteq[d]\setminus S$, consider the star configuration in which, among the nonblue-child edges, $e_i$ is open exactly when $i\in A$. 

Its unnormalized conditional weight, given $\eta$ and $\mathcal B_v^c$, is
\[
    W_\eta(A)
    =
    C(\eta)
    (1-p)^b
    p^{|A|}
    (1-p)^{d-b-|A|}
    q^{-|A|},
\]
where $C(\eta)>0$ does not depend on $A$. Indeed, each nonblue-child open edge merges one distinct finite child component into the component of $v$, decreasing the cluster count by one. After removing the factors independent of $A$, the weight is proportional to $\left(p/q\right)^{|A|}(1-p)^{d-b-|A|}$. Normalizing over $A\subseteq[d]\setminus S$ shows that the $d-b$ nonblue-child edges are independent Bernoulli variables with parameter $\widehat p$.

Let $\widehat{\mathcal F}_v^{\mathrm{ext}}:=\mathcal F_v^{\mathrm{ext}}\vee\sigma\left(I_e:e\notin E(\mathcal T_v^{\mathrm{ext}})\right)$, and let $\mathcal Q_n(v)$ be the event that the exploration started from $v$ either selects a probe edge before generation $n$ or reaches generation $n$. We claim inductively that
\begin{equation}
\label{eq:cavity-finite-depth}
    \overline\nu\left(
        \mathcal Q_n(v)
        \,\middle|\,
        \widehat{\mathcal F}_v^{\mathrm{ext}},
        \mathcal B_v^c
    \right)
    \ge\delta.
\end{equation}

The case $n=0$ is immediate. Assume~\eqref{eq:cavity-finite-depth} for $n$, and condition on
\[
    S_v:=\{w:w\text{ is a blue child of }v\}=S,
    \qquad
    |S|=b.
\]
Order the children as $w_1,\ldots,w_d$ and process their branches in this order. Before processing $w_j$, let $\mathcal G_{j-1}$ be generated by $\widehat{\mathcal F}_v^{\mathrm{ext}}$, $\mathbf 1_{\mathcal B_v}$, $S_v$, and the probe marks, edge states, and descendant-exploration outcomes revealed while processing $w_1,\ldots,w_{j-1}$. 

In the $j$th branch, first reveal the probe mark on $e_j:=\{v,w_j\}$. If $w_j$ is blue and $e_j$ is a probe, the exploration selects $e_j$. If $w_j$ is nonblue and $e_j$ is not a probe, reveal its state; when it is open, continue with the depth-$n$ exploration rooted at $w_j$. Let $\mathcal G_j$ be the sigma-field after this processing, and let $\mathcal D_j$ be the event that the $j$th branch produces neither a selected probe edge nor a continuation to depth $n$.

Fix $j$ and consider $\{S_v=S\}\cap\bigcap_{k<j}\mathcal D_k$. If $w_j\notin S$, then no information from the descendant branch $\mathcal T_{w_j}^{\mathrm{ext}}$ has yet been exposed, apart from the event $\mathcal B_{w_j}^c$. After revealing the nonprobe mark and the state of $e_j$, set $\mathcal H_j:=\mathcal G_{j-1}\vee\sigma(I_{e_j},\omega(e_j))$. The reveal order gives the sigma-field inclusion $\mathcal H_j\subseteq\widehat{\mathcal F}_{w_j}^{\mathrm{ext}}\vee\sigma(\mathbf 1_{\mathcal B_{w_j}})$.

The induction hypothesis at $w_j$, followed by the tower property from $\widehat{\mathcal F}_{w_j}^{\mathrm{ext}}\vee\sigma(\mathbf 1_{\mathcal B_{w_j}})$ down to $\mathcal H_j$, therefore gives conditional probability at least $\delta$ for $\mathcal Q_n(w_j)$ whenever $w_j$ is nonblue and $e_j$ is open. The star calculation, which applies Lemma~\ref{lem:conditional-star-factorization} after forcing every edge to a blue child to be closed on $\mathcal B_v^c$, shows that, conditional on $\mathcal B_v^c$ and $S_v=S$, the remaining nonblue-child edges are independent Bernoulli variables with parameter $\widehat p$. This remains true after conditioning on the previously processed branches. Finally, the probe marks are independent of the random-cluster configuration and of one another. Combining these three facts gives, on $\bigcap_{k<j}\mathcal D_k$,
\[
    \overline\nu\left(
        \mathcal D_j
        \,\middle|\,
        \mathcal G_{j-1}
    \right)
    \le
    \begin{cases}
        1-\varepsilon,
        & w_j\in S,\\
        1-(1-\varepsilon)\widehat p\,\delta,
        & w_j\notin S.
    \end{cases}
\]

Iterating the tower property along $(\mathcal G_j)_{j=0}^d$ gives
\[
    \overline\nu\left(
        \mathcal Q_{n+1}(v)^c
        \,\middle|\,
        \widehat{\mathcal F}_v^{\mathrm{ext}},
        \mathcal B_v^c,
        S_v=S
    \right)
    \le
    (1-\varepsilon)^b
    \left(
        1-(1-\varepsilon)\widehat p\,\delta
    \right)^{d-b}
    =
    1-\Psi_b(\delta).
\]
The factor $(1-\varepsilon)^b$ corresponds to having no probe among the blue children, while the second factor corresponds to having no successful continuation through the nonblue children. Thus, for every $S$,
\[
    \overline\nu\left(
        \mathcal Q_{n+1}(v)
        \,\middle|\,
        \widehat{\mathcal F}_v^{\mathrm{ext}},
        \mathcal B_v^c,
        S_v=S
    \right)
    \ge
    \Psi_{|S|}(\delta)
    \ge
    \delta.
\]
Averaging over $S_v$ proves~\eqref{eq:cavity-finite-depth} for $n+1$.

Here generations are measured by graph distance from the starting vertex. The events $\mathcal Q_n(u)$ decrease with $n$. If $\bigcap_{n\ge0}\mathcal Q_n(u)$ occurs without selecting an edge, then the finitely branching tree of active vertices contains vertices at every depth. K\"onig's infinity lemma therefore gives an infinite downward ray of open nonprobe edges from $u$, which implies $\mathcal B_u$. Conditional continuity from above in~\eqref{eq:cavity-finite-depth} now shows that, on $\mathcal B_u^c$, the probability that the exploration selects a probe edge is at least $\delta$.

Let $E_u$ denote the probe edge selected by the exploration, and set $\mathcal S_u:=\{E_u\text{ is defined}\}$. For a fixed oriented edge $e=\{v,w\}$, let $\mathcal A_e:=\{E_u=e\}$. To determine whether $\mathcal A_e$ occurs, the exploration checks the revealed path from $u$ to $v$, the earlier tested branches, the probe marks, and the event $\mathcal B_w$. It never checks whether $e$ itself is open or closed. Therefore $\mathcal A_e$ is measurable from the probe marks and the configuration off $e$.

Write $\mathbf I$ for the complete family of probe marks $(I_f)$, where $f$ ranges over all oriented child edges of $\mathcal T_\infty$. For every edge $e$, the single-edge conditional open probability is either $p$ or $\widehat p$. Since $\mathbf I$ is independent of the random-cluster configuration, the same is true under $\overline\nu$, and hence
\[
    \frac{
        \overline\nu\left(
            \omega(e)=1
            \,\middle|\,
            \omega_{E(\mathcal T_\infty)\setminus\{e\}},
            \mathbf I
        \right)
    }{
        \overline\nu\left(
            \omega(e)=0
            \,\middle|\,
            \omega_{E(\mathcal T_\infty)\setminus\{e\}},
            \mathbf I
        \right)
    }
    \ge
    \frac{p}{1-p}.
\]
Since $\mathcal A_e$ and $\mathcal F_u^{\mathrm{ext}}$ are measurable off $e$, integrating the preceding inequality over the configurations off $e$ gives
\[
    \overline\nu\left(
        \mathcal A_e\cap\{\omega(e)=1\}
        \,\middle|\,
        \mathcal F_u^{\mathrm{ext}}
    \right)
    \ge
    \frac{p}{1-p}
    \overline\nu\left(
        \mathcal A_e\cap\{\omega(e)=0\}
        \,\middle|\,
        \mathcal F_u^{\mathrm{ext}}
    \right).
\]
On $\mathcal A_e\cap\{\omega(e)=1\}$, the active open path from $u$ reaches the parent endpoint of $e$, and the child endpoint is blue; hence $\mathcal B_u$ occurs. On $\mathcal B_u^c\cap\mathcal S_u$, the selected edge is necessarily closed. The tree has countably many edges, so the edgewise conditional inequalities may be taken to hold on one common full-measure event. Since the events $(\mathcal A_e)_e$ are disjoint, applying the inequality first to finite collections of possible selected edges and then using conditional monotone convergence gives
\[
\begin{aligned}
    \nu\left(
        \mathcal B_u
        \,\middle|\,
        \mathcal F_u^{\mathrm{ext}}
    \right)
    &\ge
    \frac{p}{1-p}
    \overline\nu\left(
        \mathcal B_u^c\cap\mathcal S_u
        \,\middle|\,
        \mathcal F_u^{\mathrm{ext}}
    \right).
\end{aligned}
\]
Let $a_u:=\nu\left(\mathcal B_u\,\middle|\,\mathcal F_u^{\mathrm{ext}}\right)$. By the preceding exploration estimate, $\overline\nu\left(\mathcal B_u^c\cap\mathcal S_u\,\middle|\,\mathcal F_u^{\mathrm{ext}}\right)\ge\delta(1-a_u)$. Therefore
\[
    a_u
    \ge
    \frac{p}{1-p}\delta(1-a_u)
    \implies
    a_u
    \ge
    \frac{p\delta}{1-p+p\delta}
    =:
    \theta
    >0.
\]
The constant $\theta$ depends only on $p,q,$ and $d$ through $\delta$.
\end{proof}

\end{document}
